\chapter{Ответы и решения}\label{app:answers}

\newcommand{\ans}[1]{\par\smallskip\noindent\textbf{\color{hydr}#1.}\ }

\section*{Глава 1}
\ans{1.1} Если ось $c$ лежит в плоскости $r$--$\theta$ под углом $\pm\chi$ к радиусу, то
$\cos^2$ с радиусом равен $\cos^2\chi$, с окружностью --- $\sin^2\chi$, поэтому
$f_\theta/(f_r+f_\theta)=\sin^2\chi$ (часть осей выходит к оси трубы, отсюда нормировка на
$f_r+f_\theta$). Для Lepine $\sin^2\chi=0.30/0.89=0.337$, $\chi\approx35^\circ$. С учётом разброса
осей в модели берут средний наклон $\chi_0=30^\circ$ и разброс $26^\circ$.

\ans{1.2} $f=178/15\,660=1.14\%$. Площадь гидрида $0.0114\cdot240^2\approx655$~мкм$^2$, пластинка
$5\times0.6=3$~мкм$^2$ --- около 220 пластинок. В модели пластинки короче (медиана полудлины 0.8~мкм),
поэтому их 400--600.

\section*{Глава 2}
\ans{2.1} (а) $\lg C=5.56-2240/673.15=2.232$, $C\approx171$~ppm. (б) $T=2240/(5.56-\lg178)=2240/3.310=676.8$~K
$=404\,^\circ$C. (в) $T=1680/(5.08-2.250)=593.7$~K $=321\,^\circ$C. (г) $Q_{\TSSD}=\ln10\cdot1000\cdot2.24\cdot8.314=42.9$~кДж/моль,
$Q_{\TSSP}=32.2$~кДж/моль.

\ans{2.2} При 573~K: $c_{\TSSP}=10^{5.08-2.931}=141$~ppm, $c_{\TSSD}=10^{5.56-3.908}=44.9$~ppm, отношение 3.14.
Цена гистерезиса $0.84\cdot573\cdot\ln3.14\approx550$~МПа. Упругая энергия несоответствия
$E\eps^{*2}\approx90\,000\cdot0.0036\approx320$~МПа --- тот же порядок. Гистерезис --- это цена
размещения несоответствия (упругого и пластического) при зарождении; так его и объясняет Пульс.

\ans{2.3} $n_HQ/T$: при 523~K --- 6.2, при 673~K --- 4.8~МПа/$^\circ$C. В интервале выпадения «курс»
меняется на $\pm13\%$; в модели он считается при текущей температуре.

\section*{Глава 3}
\ans{3.1} $\lambda=71.4$~ГПа, $\mu=33.6$~ГПа; $\eps=350/90\,000=0.39\%$.

\ans{3.2} $d\sigma/dx=0$ --- напряжение одинаково вдоль стержня; среднее равно нулю --- значит, нулю
везде. Стержень просто удлиняется на $f\eps^*L$. В 1D любая собственная деформация
\emph{совместна}: её можно «проинтегрировать» в смещение. В 2D и 3D локальное расширение включения
несовместно с окружением: нельзя раздуть кусок, не раздвинув соседей, --- отсюда внутренние напряжения.

\ans{3.3} Ось пластинки $\mathbf t=(\cos\psi,\sin\psi)$, нормаль $\mathbf n=(-\sin\psi,\cos\psi)$ в
осях $(\theta,r)$; $\eps^*=\eps_t\mathbf t\mathbf t^{\mathsf T}+\eps_n\mathbf n\mathbf n^{\mathsf T}+\eps_t\mathbf e_z\mathbf e_z^{\mathsf T}$,
$\eps^*_{\theta\theta}=\eps_t\cos^2\psi+\eps_n\sin^2\psi$, откуда \eqref{eq:gpsi}. При 150~МПа
$\Delta g=54.6$~МПа, $\eps_n\Delta g=3.9$~МПа, это $0.73\,^\circ$C против измеренных $0.08\cdot150=12\,^\circ$C ---
в 16 раз меньше.

\ans{3.4} $\lambda f\eps_t=71\,360\cdot0.0114\cdot0.0458=37$~МПа, $(\lambda+2\mu)f\eps_t=72$~МПа;
$g=(0.0458\cdot(-37)+0.072\cdot(-37)+0.0458\cdot(-72))/0.072\approx-107$~МПа, в движущей силе
$-7.7$~МПа $\approx-1.4\,^\circ$C. Совпадает со сверкой с МКЭ ($-94$~МПа на 1\%).

\section*{Глава 4}
\ans{4.1} $s=\sigma_1\,\mathrm{diag}(2/3,-1/3,-1/3)$, $s:s=\tfrac23\sigma_1^2$, $\sigma_{\text{экв}}=|\sigma_1|$.
При $\sigma=-pI$ девиатор нулевой, $\sigma_{\text{экв}}=0$: всестороннее сжатие текучести не вызывает,
поэтому металл перед кончиком держит $\sigma_h\approx-500$~МПа. Течёт он от девиаторной части.

\ans{4.2} Если $\eps_t$ в плоскости снято течением, окружная пластинка под $\sigma_\theta$ работы не
получает: $g_{\text{окр}}\approx0$. У радиальной $g_{\text{рад}}=\sigma_\theta$. Разница $\approx\sigma_\theta$ ---
верхняя оценка, расчёт даёт $1.0\,\sigma_\theta$.

\ans{4.3} В перемычке между пластинками зоны течения перекрываются: металл уже течёт от первой
пластинки, и вторая его «досгибает» с другим, упрочнённым пределом; сумма двух решений не
удовлетворяет условию текучести. Под нагрузкой вдоль окружной пластинки течение облегчено
нагрузкой с одной стороны и затруднено с другой --- нелинейность сильнее, отсюда большая ошибка.

\section*{Глава 5}
\ans{5.1} $d\Delta G/dr=8\pi r\gamma-4\pi r^2\Delta=0\Rightarrow r^*=2\gamma/\Delta$;
$\Delta G^*=16\pi\gamma^3/(3\Delta^2)$. При $+10\%$ движущей силы барьер в $1.21$ раза ниже ($-17\%$).

\ans{5.2} $d\ln r/d\Delta=2B\Delta_0^2/\Delta^3$, при $\Delta_0$ --- $2B/\Delta_0=10.7$ на МПа, $\times5.4=58$
на градус. Разница $0.2\,^\circ$C --- множитель $e^{11.5}\approx10^5$.

\ans{5.3} $f=(2.39)(0.61)^2/4=0.222$, $\sqrt f=0.47$, сдвиг $15\cdot0.53=7.9$~МПа $=1.47\,^\circ$C.
При $\gamma_{\alpha\alpha}=0.18$: $\cos\theta=0.51$, $f=0.15$, сдвиг 9.1~МПа $=1.7\,^\circ$C. Фора границы
слабо чувствительна к энергиям --- хорошо для обоснованности.

\ans{5.4} $12.4/0.08=155$~МПа; $(5+1.5)/0.08=81$~МПа. Оба --- как в опыте.

\ans{5.5} Поле и водород пересчитываются между шагами. Если за шаг рождается много зародышей, они
не видят полей друг друга и берут один и тот же водород: пластинок и гидрида получается больше,
чем водорода (такая ошибка в модели была и исправлена), а цепочки --- случайнее.

\section*{Глава 6}
\ans{6.1} $D(300\,^\circ\text{C})=7.1\cdot10^{-11}$~м$^2$/с, за минуту 92~мкм; $D(200\,^\circ\text{C})=9.9\cdot10^{-12}$,
за минуту 34~мкм. Много больше зерна: водород выравнивается по многим зёрнам; шаг пакетов
57~мкм --- того же порядка, что диффузионная длина: каждый пакет «осушает» свою окрестность.

\ans{6.2} $3\cdot15\,660\approx4.7\cdot10^4$~ppm$\cdot$мкм$^2$; при 20~ppm --- площадь 2350~мкм$^2$,
радиус $\approx27$~мкм.

\ans{6.3} $D/h=118$~мкм/с, $v=0.05\cdot118\cdot20/15\,660=7.6\cdot10^{-3}$~мкм/с $\approx0.45$~мкм/мин.
От 1.2 до 5~мкм (по 1.9~мкм с каждого конца) --- около 4~мин. При 15\,$^\circ$C/мин за это время металл
остынет на 60\,$^\circ$C, при 0.3\,$^\circ$C/мин --- на 1.3\,$^\circ$C: рост во времени --- кандидат на
объяснение влияния скорости охлаждения.

\ans{6.4} $V_H\sigma_h/RT=1.7\cdot10^{-6}\cdot10^8/(8.314\cdot603)=0.034$: $+3.4\%$. При $-540$~МПа:
$e^{-0.18}=0.83$, на 17\% меньше.

\section*{Глава 7}
\ans{7.1} $240^2/(2.5\cdot4.5)=5120$ зёрен. Крутой след --- отклонение оси $c$ от радиуса больше
$65^\circ$: при $\chi\sim N(\pm30^\circ,26^\circ)$ это $P(z>1.35)\approx9\%$ зёрен.

\ans{7.2} $(415-100)/0.5=630$ шагов, 105~мин модельного времени. При 0.3\,$^\circ$C/мин шагов столько же,
но каждый в 10 раз длиннее по времени (1050~мин): диффузия за шаг сильнее, а время счёта почти то же.

\section*{Глава 8}
\ans{8.1} Числитель $6\cdot0.5+8+2=13$, знаменатель 30: RHF $=0.43$. Если 6-микронный след --- три
окружные пластинки, «по пластинкам» он даёт вес 0: RHF $=10/30=0.33$. Стопка видна на шлифе как
наклонный гидрид, а пластинки в ней окружные --- меры отвечают на разные вопросы.

\ans{8.2} Цена $400+200/50=404$, RHCP $=(1-404/600)/0.98=0.33$. Для 0.5 нужно
$600-0.98x=306$, $x=300$~мкм по гидриду.

\ans{8.3} По формуле Коши средняя хорда выпуклой фигуры для изотропных случайных прямых ---
$\pi A/P$; для круга $\pi\cdot\pi R^2/(2\pi R)=\pi R/2$. Изотропная плоскость, задевшая диск, высекает
в его плоскости изотропную прямую.

\section*{Глава 9}
\ans{9.1} Водород определяет, \emph{при какой температуре} начнётся выпадение, а порог --- \emph{какое
семейство} начнёт первым; это зависит от разницы фор в градусах ($b$ против $0.08\,\sigma_\theta$),
а не от температуры. При 60 ppm выпадение начинается при низкой температуре, диффузия медленнее,
водород не успевает стечься к «победителям», и зарождаются все семейства --- в том числе крутые
зёрна без нагрузки.

\ans{9.2} Пластинки --- около базиса своих зёрен, а стопка из них --- ступенчатая и идёт по радиусу.
Снимок видит стопку. Для сравнения с Lepine правильна мера по снимку; значит, «радиальность»
гидрида на шлифе --- свойство стопки, а не отдельной пластинки.

\section*{Глава 10}
\ans{10.1} Радиальная ($\mathbf n=\mathbf e_\theta$): $\sigma:\eps^*=\sigma_\theta\eps_n+\sigma_z\eps_t$; окружная
($\mathbf n=\mathbf e_r$): $\sigma_\theta\eps_t+\sigma_z\eps_t$. Разность $(\eps_n-\eps_t)\sigma_\theta$ от
$\sigma_z$ не зависит. В матрице Мизеса течение определяет девиатор, а выбор между двумя
ориентациями, переходящими друг в друга поворотом на $90^\circ$ вокруг оси трубы, зависит только
от $\sigma_\theta-\sigma_r$; $\sigma_z$ входит в обе одинаково.

\ans{10.2} Отжиг для снятия напряжений без рекристаллизации снимает остаточные напряжения, но
почти не трогает дислокационную структуру. Если фора (и порог) после такого отжига падают ---
причина в напряжениях; если нет --- в дислокациях. Порог измеряется как обычно, по RHF после
охлаждения под нагрузкой; дополнительно --- нейтронография остаточных напряжений по семействам зёрен.
